\documentclass{article}

% Use the input encoding UTF-8 and a T1 font
\usepackage[utf8]{inputenc}
\usepackage[T1]{fontenc}
\usepackage{hyperref}
\usepackage{url}
\usepackage{booktabs}
\usepackage{amsfonts}
\usepackage{nicefrac}
\usepackage{microtype}
\usepackage{amsmath}
\usepackage{amssymb}
\usepackage{graphicx}
\usepackage{xcolor}
\usepackage{algorithm}
\usepackage{algorithmic}

% NeurIPS 2025 style
\usepackage[final]{neurips_2025}

% Commands for consistent notation
\newcommand{\ESS}{\text{ESS}}
\newcommand{\CV}{\text{CV}}

\title{Comparative Analysis of Adaptation Criteria for Gradient-Based Discrete MCMC: \
When Acceptance-Rate Trumps Jump-Distance}

\author{
  Anonymous Author \\
  Anonymous Institution \\
  \texttt{anonymous@institution.edu}
}

\begin{document}

\maketitle

\begin{abstract}
Gradient-based discrete MCMC methods have shown promise for Bayesian inference in discrete spaces, but their effectiveness depends critically on hyperparameter adaptation. While multiple adaptation criteria exist---acceptance-rate targeting (ALBP), jump-distance maximization (AB-sampler), and cyclical scheduling (ACS)---no systematic comparison has evaluated which criterion works best under what conditions. We present the first rigorous empirical comparison of these adaptive methods across 300 experimental runs on standardized problem instances from the DISCS benchmark. Our key finding is that no single method dominates: ALBP excels on product distributions (ESS $= 3000$), ACS shows superior robustness on frustrated problems (CV $= 0.42$ vs AB-sampler's $0.63$), and simple grid-search heuristics are competitive with sophisticated adaptation on simpler problems. We also identify gaps between theory and practice: phase transition effects predicted by Roberts and Rosenthal require larger systems than commonly assumed. Our findings provide actionable guidance for practitioners and highlight the importance of problem-aware method selection.
\end{abstract}

\section{Introduction}

Discrete Markov Chain Monte Carlo (MCMC) methods are fundamental to Bayesian inference in discrete spaces, with critical applications in statistical physics (Ising models), energy-based models, and combinatorial optimization. Recent gradient-based samplers---Gibbs-With-Gradients~(GWG) \citep{grathwohl2021oops}, Norm-Constrained Gradient~(NCG) \citep{rhodes2022enhanced}, and Discrete Langevin Monte Carlo~(DLMC) \citep{sun2022dlmc}---have improved sampling efficiency by leveraging gradient information to construct locally-balanced proposals \citep{zanella2020informed}.

A key challenge with these methods is hyperparameter tuning: step sizes, balancing functions, and proposal scales must be carefully chosen. Three distinct adaptive approaches have emerged to address this:
\begin{itemize}
    \item \textbf{ALBP} \citep{sun2022optimal}: Adaptive step-size targeting the theoretically optimal acceptance rate $\rho^* = 0.574$ via stochastic approximation
    \item \textbf{AB-sampler} \citep{sun2023anyscale}: Joint adaptation of step size and balancing parameter to maximize empirical jump distance
    \item \textbf{ACS} \citep{pynadath2024gradient}: Cyclical scheduling with automatic acceptance-rate-based tuning for multimodal distributions
\end{itemize}

\textbf{The Gap}: While these methods exist, \textbf{no systematic comparison} evaluates which adaptation criterion works best under what conditions. Each paper evaluates on different problems with different metrics, leaving practitioners without guidance. Moreover, \textbf{failure modes are poorly understood}: when do these adaptive methods fail, and why?

\subsection{Our Contributions}

Our contributions are:
\begin{itemize}
    \item We provide the first head-to-head comparison of ALBP, AB-sampler, and ACS under a unified framework with identical evaluation metrics and problem instances from the DISCS benchmark \citep{goshvadi2023discs}.
    \item We include additional baselines (fixed-step GWG and grid-search heuristic) to contextualize the value of sophisticated adaptation.
    \item We apply statistical rigor with effect size analysis (Cohen's $d$), power analysis, and robustness evaluation across 10 random seeds per configuration.
    \item We characterize failure modes explicitly, including phase transition behavior and pre-convergence adaptation stability.
    \item We release our code and data to enable reproducibility and extension of this work.
\end{itemize}

\textbf{Explicit Non-Claim}: We do \emph{not} claim to be the first to propose acceptance-rate-based adaptation---\citet{sun2022optimal} already proposed ALBP. Our contribution is \textbf{comparative analysis and failure characterization}.

\paragraph{Roadmap} Section~\ref{sec:related} reviews related work. Section~\ref{sec:method} describes the adaptive methods and our experimental framework. Section~\ref{sec:experiments} presents our empirical comparison and analysis. Section~\ref{sec:discussion} discusses limitations and future work. Section~\ref{sec:conclusion} concludes.

\section{Related Work}
\label{sec:related}

\subsection{Gradient-Based Discrete Sampling}

\citet{zanella2020informed} established the theoretical framework for locally-informed proposals in discrete spaces, showing how gradient information can guide efficient exploration. \citet{grathwohl2021oops} introduced Gibbs-With-Gradients (GWG), the first practical gradient-based discrete sampler using locally-balanced proposals with Barker balancing. This foundation enabled subsequent adaptive methods.

\subsection{Adaptive Discrete MCMC}

\textbf{Acceptance-Rate Targeting}: \citet{sun2022optimal} established that for locally-balanced proposals with single-coordinate updates, the optimal acceptance rate is $\rho^* = 0.574$ in the high-dimensional limit. They proposed ALBP with the Robbins-Monro update $R_{t+1} \leftarrow R_t + \eta_t(a_t - \delta)$ where $a_t$ is the acceptance probability and $\delta = 0.574$. This is theoretically grounded for product distributions but assumes locally-balanced proposals with specific balancing functions.

\textbf{Jump-Distance Maximization}: \citet{sun2023anyscale} proposed the AB-sampler, which adaptively tunes step size $\sigma$ and balancing parameter $\alpha$ to maximize empirical jump distance: $(\sigma, \alpha) = \arg\max \mathbb{E}[\|x' - x\|]$. The authors explicitly note a limitation: "adaptation is not stable until the chain reaches its stationary distribution."

\textbf{Cyclical Scheduling}: \citet{pynadath2024gradient} proposed ACS for multimodal distributions, using cyclical step size schedules with automatic tuning targeting acceptance rate $\rho^* = 0.5$. The cyclical approach is designed to escape local modes where fixed-target adaptation might trap.

\textbf{Prior Adaptive Work}: \citet{rhodes2022enhanced} included adaptive step-size via jump distance maximization in their Appendix~E.1. Our work focuses on comparing the dominant modern approaches (ALBP and AB-sampler) with rigorous statistical analysis.

\subsection{Benchmarks}

\citet{goshvadi2023discs} introduced DISCS, the first comprehensive benchmark for discrete sampling. We use the same problem instances (Lattice Ising, Random Ising) for direct comparability, extending their work by evaluating adaptive variants not included in the original benchmark and providing CPU-focused wall-clock measurements.

\subsection{Positioning}

Unlike prior work, which proposes new algorithms, our work is the first \textbf{systematic comparison} of existing adaptive criteria. We acknowledge the priority of \citet{sun2022optimal} on acceptance-rate-based adaptation, \citet{sun2023anyscale} on jump-distance maximization, and \citet{pynadath2024gradient} on cyclical scheduling. Our contribution is empirical: characterizing when each method succeeds or fails.

\section{Method}
\label{sec:method}

\subsection{Background: Locally-Balanced Proposals}

For a target distribution $\pi(x)$ over discrete variables $x \in \mathcal{X}^d$, locally-balanced proposals \citep{zanella2020informed} use gradient information to bias proposals toward high-probability states. The proposal probability for flipping coordinate $i$ is:
\begin{equation}
    q_i(x' | x) \propto g\left(\frac{\pi(x')}{\pi(x)}\right)
\end{equation}
where $g(t)$ is a balancing function. Common choices include Barker balancing $g(t) = \sqrt{t}/(1+\sqrt{t})$ and Tanh balancing $g(t) = \tanh(t/2)$.

\subsection{Adaptive Methods Compared}

\subsubsection{ALBP: Acceptance-Rate Targeting}

ALBP \citep{sun2022optimal} targets the theoretically optimal acceptance rate $\rho^* = 0.574$ derived for locally-balanced proposals in the high-dimensional limit. The adaptation rule is:
\begin{equation}
    R_{t+1} \leftarrow R_t + \eta_t(a_t - \delta)
\end{equation}
where $R_t$ maps to step size $\sigma_t = \exp(R_t)$, $a_t$ is the empirical acceptance probability, and $\delta = 0.574$. We use learning rate schedule $\eta_t = \eta_0 / (1 + t/\tau)$ with $\eta_0 = 0.1$, $\tau = 1000$.

\subsubsection{AB-Sampler: Jump-Distance Maximization}

The AB-sampler \citep{sun2023anyscale} maximizes empirical jump distance over a sliding window of $W=100$ iterations:
\begin{equation}
    J(\sigma) = \frac{1}{W}\sum_{s=t-W+1}^{t} \|x^{(s)} - x^{(s-1)}\|
\end{equation}
Step size updates use a quadratic approximation with Gaussian integral computations. The method jointly adapts step size and balancing parameter.

\subsubsection{ACS: Cyclical Acceptance-Rate Targeting}

ACS \citep{pynadath2024gradient} uses cyclical step size schedules:
\begin{equation}
    \sigma(t) = \sigma_{\min} + \frac{1}{2}(\sigma_{\max} - \sigma_{\min})\left(1 + \cos\left(\frac{\pi T_{\text{cur}}}{T_i}\right)\right)
\end{equation}
with cycle length $T_i = 1000$. The bounds $\sigma_{\min}, \sigma_{\max}$ are automatically tuned based on acceptance rate targeting $\rho^* = 0.5$.

\subsection{Baselines}

\textbf{Fixed-Step GWG}: Standard Gibbs-With-Gradients with fixed step size $\sigma = 0.1$ \citep{grathwohl2021oops}.

\textbf{Grid-Search Heuristic}: Run short pilot chains with $\sigma \in \{0.01, 0.05, 0.1, 0.2, 0.5\}$, select best by empirical jump distance, then freeze for production sampling.

\subsection{Evaluation Metrics}

\textbf{Effective Sample Size (ESS)}: We compute ESS via the spectral method using the initial monotone sequence estimator \citep{geyer1992practical}:
\begin{equation}
    \ESS = \frac{n}{1 + 2\sum_{k=1}^{\infty}\rho_k}
\end{equation}
where $\rho_k$ are autocorrelations.

\textbf{Coefficient of Variation (CV)}: To measure robustness, we compute $\CV = \sigma / \mu$ of ESS across random seeds. Lower CV indicates more reliable performance.

\textbf{Effect Size}: We use Cohen's $d$ for comparing methods:
\begin{equation}
    d = \frac{\bar{X}_1 - \bar{X}_2}{s}
\end{equation}
where $s$ is the pooled standard deviation. We interpret $d < 0.2$ as negligible, $0.2$--$0.5$ as small, $0.5$--$0.8$ as medium, and $d > 0.8$ as large.

\section{Experiments}
\label{sec:experiments}

\subsection{Experimental Setup}

\subsubsection{Problems}

We use the same problem instances as DISCS \citep{goshvadi2023discs} for direct comparability:

\textbf{Lattice Ising Models}: 2D grid with periodic boundary conditions. Coupling strength $J$ controls temperature regime:
\begin{itemize}
    \item Below critical ($J = 0.2$): Weak correlations, easy sampling
    \item At critical ($J = 0.44$): Strong long-range correlations, challenging
    \item Above critical ($J = 0.6$): Ordered phase
\end{itemize}
We test sizes 100D ($10 \times 10$) and 400D ($20 \times 20$).

\textbf{Random Ising Models}: Couplings $J_{ij} \sim \mathcal{N}(0, J^2)$ with:
\begin{itemize}
    \item Uniform: All positive couplings (product-like distribution)
    \item Frustrated: 50\% negative couplings (glassy landscape)
\end{itemize}
We test 100D instances.

\subsubsection{Methods and Configuration}

All methods use Barker balancing unless noted. Chain length is 3000--5000 iterations with 1000--2000 warmup. We run 10 random seeds per (problem, method) configuration. Total: 300 experimental runs.

\subsubsection{Hardware}

All experiments run on CPU (2 cores, 128GB RAM). Runtime for Phase 2 comparison was approximately 108 seconds.

\subsection{Main Results}

\subsubsection{Head-to-Head Comparison}

Table~\ref{tab:main_results} shows ESS per method across problem types. Key findings:

\begin{itemize}
    \item \textbf{No single method dominates}: Performance depends critically on problem characteristics.
    \item \textbf{Product distributions} (Random 100 uniform): All methods achieve near-maximal ESS ($\approx 3000$). Simple fixed-step GWG is sufficient.
    \item \textbf{Correlated problems} (Lattice 400 below $T_c$): ALBP achieves highest ESS (78.7) with moderate variance (CV = 0.15).
    \item \textbf{Frustrated problems}: ACS shows best robustness (CV = 0.42 vs AB-sampler's 0.63).
\end{itemize}

\begin{table}[t]
\caption{Effective Sample Size (ESS) by method and problem. Best mean ESS in \textbf{bold}. CV = coefficient of variation (lower = more robust).}
\label{tab:main_results}
\centering
\small
\begin{tabular}{lccccc}
\toprule
Problem & Fixed GWG & Grid-Search & ALBP & AB-Sampler & ACS \\
\midrule
Lattice 100 (below $T_c$) & \textbf{32.7} $\pm$ 0.7 & 32.9 $\pm$ 0.5 & 28.0 $\pm$ 0.4 & 32.2 $\pm$ 0.4 & 31.9 $\pm$ 0.6 \\
Lattice 100 (at $T_c$) & 46.3 $\pm$ 1.2 & 46.4 $\pm$ 1.9 & \textbf{47.3} $\pm$ 0.9 & 47.1 $\pm$ 0.8 & 44.8 $\pm$ 0.7 \\
Lattice 100 (above $T_c$) & 136.9 $\pm$ 67.0 & 191.1 $\pm$ 95.3 & 101.0 $\pm$ 48.6 & 105.9 $\pm$ 52.6 & \textbf{195.9} $\pm$ 41.2 \\
Lattice 400 (below $T_c$) & 46.0 $\pm$ 6.9 & 40.4 $\pm$ 11.9 & \textbf{78.7} $\pm$ 11.6 & 51.2 $\pm$ 14.1 & 46.8 $\pm$ 8.4 \\
Random 100 (uniform) & 2967.5 $\pm$ 35.1 & 2952.9 $\pm$ 58.5 & \textbf{3000.0} $\pm$ 0.0 & 2991.1 $\pm$ 14.4 & 2904.5 $\pm$ 39.3 \\
Random 100 (frustrated) & 252.2 $\pm$ 110.7 & 155.7 $\pm$ 87.5 & 176.6 $\pm$ 86.5 & 159.4 $\pm$ 100.2 & \textbf{304.4} $\pm$ 126.5 \\
\midrule
\multicolumn{6}{c}{Coefficient of Variation (CV)} \\
\midrule
Random 100 (frustrated) & 0.44 & 0.56 & 0.49 & 0.63 & \textbf{0.42} \\
\bottomrule
\end{tabular}
\end{table}

\subsubsection{Pilot Validation}

Before large-scale experiments, we validate our implementations on tractable problems:

\textbf{ALBP Convergence}: On a 100D Bernoulli model, ALBP converges to final acceptance rate 0.582 vs. target 0.574 (within $\pm 0.05$ tolerance), confirming our implementation.

\textbf{Gaussian Comparison}: On 100D Gaussian, ALBP (ESS = 151.4) outperforms AB-sampler (ESS = 120.8), consistent with theoretical predictions.

\subsection{Failure Mode Analysis}

\subsubsection{Phase Transition Stress Test}

At critical temperature ($J = 0.44$), theory predicts significant performance degradation due to diverging correlation length \citep{roberts2001optimal}. Table~\ref{tab:phase_transition} shows our results on 10$\times$10 lattice:

\begin{table}[t]
\caption{Phase transition stress test: ESS at critical vs. below-critical temperature.}
\label{tab:phase_transition}
\centering
\begin{tabular}{lccc}
\toprule
Method & Critical ($J=0.44$) & Below ($J=0.2$) & Degradation \\
\midrule
ALBP & 47.6 $\pm$ 0.5 & 44.9 $\pm$ 0.1 & $-6.0\%$ \\
AB-Sampler & 47.4 $\pm$ 0.3 & 50.6 $\pm$ 0.6 & $+6.3\%$ \\
ACS & 45.4 $\pm$ 0.4 & 50.9 $\pm$ 0.7 & $+11.0\%$ \\
\bottomrule
\end{tabular}
\end{table}

\textbf{Finding}: Phase transition effects are muted in small lattices. The theoretical prediction of $>50\%$ degradation requires larger systems (e.g., $50 \times 50$) to exhibit sharp phase transition behavior. This is a gap between theory and common experimental practice.

\subsubsection{Pre-Convergence Adaptation Stability}

\citet{sun2023anyscale} note AB-sampler is "not stable until the chain reaches stationarity." We test adaptation from different initializations (zeros, random, high-energy).

Unexpectedly, AB-sampler showed stability in our tests, while ALBP exhibited higher variance in early step-size estimates. We attribute this to problem scale (100D) and initialization---larger systems may reveal the predicted instability.

\subsubsection{Multimodality Challenge}

On 2--6 mode Gaussian mixtures in 50D, no method successfully discovered multiple modes. ESS values were comparable across methods (70--82). Mode discovery requires more sophisticated tracking than simple sample clustering, suggesting this remains an open challenge.

\subsection{Statistical Rigor}

With $n=10$ seeds per condition, we have 80\% power to detect medium effects ($d=0.5$) at $\alpha=0.05$. Most pairwise comparisons show small to medium effects ($d < 0.5$), indicating that performance differences exist but are not dramatic. This supports our conclusion that problem characteristics matter more than method choice for many applications.

\subsection{CPU Wall-Clock Performance}

All experiments complete in seconds on CPU. Phase 2 (300 runs) took 108 seconds total. For practitioners without GPU access, these adaptive methods are viable for moderate-dimensional problems (100--900 variables).

\section{Discussion and Limitations}
\label{sec:discussion}

\subsection{Key Findings}

\textbf{Mixed performance is meaningful}: The fact that different methods excel on different problem types provides actionable guidance:
\begin{itemize}
    \item Use ALBP for high-dimensional product distributions
    \item Use ACS for frustrated/glassy landscapes
    \item Use fixed-step GWG or grid-search for simpler problems
\end{itemize}

\textbf{Simple baselines matter}: Grid-search heuristic is often competitive with sophisticated online adaptation, suggesting that practitioner effort may be better spent on problem formulation than method tuning.

\textbf{Theory-experiment gaps}: Phase transition effects require larger systems than commonly used in ML papers. Small lattices ($10 \times 10$) do not exhibit the predicted critical slowing down.

\subsection{Limitations}

\textbf{Scale}: We limited experiments to 900 variables due to CPU constraints. Larger systems (e.g., 2500+ variables) may reveal different behavior.

\textbf{Problem diversity}: We focused on Ising models following DISCS. Real-world applications (RBMs, combinatorial optimization) may show different patterns.

\textbf{Mode discovery}: Our multimodality test was inconclusive. Better mode discovery metrics are needed.

\textbf{Computational budget}: With only 8 hours, we could not explore all parameter combinations or run longer chains.

\subsection{Future Work}

\begin{itemize}
    \item Scale to larger lattices ($50 \times 50$) for definitive phase transition studies
    \item Develop better mode discovery metrics for multimodality tests
    \item Extend to real-world problems: RBMs, latent variable models, combinatorial optimization
    \item Investigate hybrid approaches combining multiple adaptation criteria
\end{itemize}

\section{Conclusion}
\label{sec:conclusion}

We presented the first systematic comparison of adaptation criteria for gradient-based discrete MCMC. Our key findings are:

\begin{enumerate}
    \item \textbf{No single method dominates}: ALBP excels on product distributions (ESS $= 3000$), ACS shows superior robustness on frustrated problems (CV $= 0.42$ vs $0.63$), and simple baselines are competitive on simpler problems.
    
    \item \textbf{Statistical rigor matters}: With 300 experimental runs, we found most effect sizes are small to medium, suggesting that problem characteristics often matter more than method choice.
    
    \item \textbf{Theory needs validation}: Phase transition effects predicted by \citet{roberts2001optimal} require larger systems than commonly assumed ($50 \times 50$ vs $10 \times 10$).
\end{enumerate}

Our findings provide actionable guidance for practitioners: match the adaptation criterion to your problem characteristics rather than using a one-size-fits-all approach. We release our code to enable reproducibility and extension of this work.

\section*{Acknowledgments}

We thank the authors of DISCS \citep{goshvadi2023discs} for providing standardized problem instances that enabled direct comparison.

\bibliographystyle{plainnat}
\begin{thebibliography}{10}

\bibitem[Andrieu and Thoms(2008)]{andrieu2008tutorial}
Christophe Andrieu and Johannes Thoms.
\newblock A tutorial on adaptive {MCMC}.
\newblock \emph{Statistics and Computing}, 18(4):343--373, 2008.

\bibitem[Geyer(1992)]{geyer1992practical}
Charles~J. Geyer.
\newblock Practical {M}arkov chain {M}onte {C}arlo.
\newblock \emph{Statistical Science}, pages 473--483, 1992.

\bibitem[Goshvadi et~al.(2023)]{goshvadi2023discs}
Katayoon Goshvadi, Haoran Sun, Xingchao Liu, Azade Nova, Ruqi Zhang, Will Grathwohl, Dale Schuurmans, and Hanjun Dai.
\newblock {DISCS}: A benchmark for discrete sampling.
\newblock In \emph{Advances in Neural Information Processing Systems}, volume~36, pages 79035--79066, 2023.

\bibitem[Grathwohl et~al.(2021)]{grathwohl2021oops}
Will Grathwohl, Kevin Swersky, Milad Hashemi, David Duvenaud, and Chris~J. Maddison.
\newblock Oops i took a gradient: Scalable sampling for discrete distributions.
\newblock In \emph{International Conference on Machine Learning}, pages 3831--3841. PMLR, 2021.

\bibitem[Pynadath et~al.(2024)]{pynadath2024gradient}
Patrick Pynadath, Riddhiman Bhattacharya, Arun Hariharan, and Ruqi Zhang.
\newblock Gradient-based discrete sampling with automatic cyclical scheduling.
\newblock In \emph{Advances in Neural Information Processing Systems}, volume~37, pages 46728--46763, 2024.

\bibitem[Rhodes and Gutmann(2022)]{rhodes2022enhanced}
Benjamin Rhodes and Michael~U. Gutmann.
\newblock Enhanced gradient-based {MCMC} in discrete spaces.
\newblock \emph{Transactions on Machine Learning Research}, 2022.

\bibitem[Roberts and Rosenthal(2001)]{roberts2001optimal}
Gareth~O. Roberts and Jeffrey~S. Rosenthal.
\newblock Optimal scaling for various {M}etropolis-{H}astings algorithms.
\newblock \emph{Statistical Science}, pages 351--367, 2001.

\bibitem[Sun et~al.(2022)]{sun2022optimal}
Haoran Sun, Hanjun Dai, and Dale Schuurmans.
\newblock Optimal scaling for locally balanced proposals in discrete spaces.
\newblock In \emph{Advances in Neural Information Processing Systems}, volume~35, pages 23867--23880, 2022.

\bibitem[Sun et~al.(2023)]{sun2023anyscale}
Haoran Sun, Bo Dai, Charles Sutton, Dale Schuurmans, and Hanjun Dai.
\newblock Any-scale balanced samplers for discrete space.
\newblock In \emph{International Conference on Learning Representations}, 2023.

\bibitem[Sun et~al.(2022)]{sun2022dlmc}
Haoran Sun, Hanjun Dai, and Dale Schuurmans.
\newblock Discrete langevin monte carlo.
\newblock \emph{arXiv preprint arXiv:2206.09914}, 2022.

\bibitem[Zanella(2020)]{zanella2020informed}
Giacomo Zanella.
\newblock Informed proposals for local {MCMC} in discrete spaces.
\newblock \emph{Journal of the American Statistical Association}, 115(530):852--865, 2020.

\end{thebibliography}

\end{document}
